\begin{abstract}%
We study how the gap between latent and intrinsic dimensionality shapes memorization in diffusion models. Evaluating the random-feature spectral theory of Bonnaire et al. at the two-atom covariance of data embedded in a larger latent space, we find that the $\dlat-\dint$ excess dimensions form a distinct noise-dimension block in the score network's feature-correlation spectrum, sitting between the $\dint$ signal modes and the $\nsamp-\dlat$ sample-specific modes. Under gradient flow every mode decays in parallel with rate equal to its eigenvalue, so the delay between generalization and memorization is the ratio of the smallest signal eigenvalue to the largest sample-specific eigenvalue, not a count of intervening modes. We show that this ratio grows with $\dlat$ because extra latent coordinates lower the pre-activation variance of the random features, de-saturating the nonlinearity and shrinking the sample-specific feature mass by more than an order of magnitude. Random-feature spectra and trainable-MLP memorization curves on synthetic data isolate this mechanism; an exactly solvable linear score model and a bridge to the closed-form empirical score connect the spectral clocks to the reported memorized fraction. On CelebA and CIFAR-10 latent diffusion, once the VAE is wide enough to preserve pixel-space sample identity, increasing latent width reduces the memorization rate. FID follows a separate quality tradeoff, improving as the VAE bottleneck is relieved and then worsening when the latent score problem becomes too wide. Finally, from the frozen VAE latent spectrum for a chosen $\dlat$, we construct a mode-absorption clock that estimates memorization percentage as a function of training step; the theory identifies the sample-bulk term this clock was missing and reduces its free parameters to one.
\end{abstract}
